\subsection{Causal reduction of the bilateral update}
\label{lib04:reduccion-causal}

The preceding unitary update realizes the transport of the state and its
memory in a common space. Observation of the first component admits an
equivalent description that explicitly retains the contribution of the
record. This description follows from the transport already constructed,
with the same coefficients and orientation: memory removed from the list
of variables reappears as dependence on earlier states.

Fix the operators \(A,D\) of \eqref{lib04:ad} and a realization
\(H\oplus H\). For \(n\geq0\), let
\begin{equation}
 v_{n+1}=Av_n-D^*m_n,\qquad
 m_{n+1}=Dv_n+A^*m_n.
 \label{lib04:recurrencia-memoria}
\end{equation}
The joint update is \(\mathscr R_a\) and preserves
\(\|v_n\|^2+\|m_n\|^2\).

\begin{proposition}[Exact elimination of the auxiliary record]
\label{lib04:prop-memoria-reducida}
For every \(n\geq0\), interpreting the empty sum as zero,
\begin{align}
 m_n&=(A^*)^n m_0+
       \sum_{k=0}^{n-1}(A^*)^{n-1-k}Dv_k,
       \label{lib04:memoria-explicita}\\
 v_{n+1}&=Av_n-D^*(A^*)^n m_0
       +\sum_{k=0}^{n-1}\mathscr K_{n-1-k}v_k,
       \label{lib04:evolucion-reducida}\\
 \mathscr K_j&=-D^*(A^*)^jD\quad(j\geq0).
       \label{lib04:nucleo-retardo}
\end{align}
The reduced equation, together with \eqref{lib04:memoria-explicita}, is
equivalent to the joint recurrence for the same data \((v_0,m_0)\).
\end{proposition}
\begin{proof}
The first identity holds for \(n=0\). If it holds for \(n\), the second
equation in \eqref{lib04:recurrencia-memoria} gives
\[
 m_{n+1}=(A^*)^{n+1}m_0+
 \sum_{k=0}^{n-1}(A^*)^{n-k}Dv_k+Dv_n,
\]
which is the identity at \(n+1\). Substitution at depth \(n\) into the
first equation proves \eqref{lib04:evolucion-reducida}. Conversely, if
\(v_n\) satisfies the reduced equation and \(m_n\) is defined by
\eqref{lib04:memoria-explicita}, the same difference of sums proves
\(m_{n+1}=Dv_n+A^*m_n\); substitution recovers the other equation.
\end{proof}

The kernel \(\mathscr K_j\) describes a transfer into the record,
\(j\) updates within it, and re-entry into the observed component.
The term \(-D^*(A^*)^nm_0\) preserves the initial memory. Omitting it
corresponds to the particular case \(m_0=0\), rather than an identity of
the general transport. Equation \eqref{lib04:evolucion-reducida} involves
only the current and earlier states. The causal character of this reduction
follows from the order of the joint recurrence.

The normalization of the balance also provides an explicit estimate.
Since \(Q_-=aI-U\), we have
\(\|A\|^2\leq r_a=(a+1)^3/((a+1)^3+a^3)<1\), whereas
\(D^*D\leq I\). Consequently,
\begin{equation}
 \|\mathscr K_j\|\leq r_a^{j/2},\qquad
 \sum_{j\geq0}\|\mathscr K_j\|
 \leq\frac{1}{1-\sqrt{r_a}}.
 \label{lib04:cota-nucleo-retardo}
\end{equation}
This bound controls the weight of memory in the reduced equation; the norm
of the joint state remains conserved by the unitary update.

The proposition is the temporal realization of the Schur reduction in
\cref{mem:prop-schur}. For the formal series
\(V(z)=\sum_{n\geq0}v_nz^n\) and
\(M(z)=\sum_{n\geq0}m_nz^n\), the recurrences yield
\(V=v_0+zAV-zD^*M\) and \(M=m_0+zDV+zA^*M\). Eliminating \(M\),
\begin{equation}
 V(z)=\bigl[I-zA+z^2D^*(I-zA^*)^{-1}D\bigr]^{-1}
       \bigl[v_0-zD^*(I-zA^*)^{-1}m_0\bigr].
 \label{lib04:schur-bilateral}
\end{equation}
The formal inverses exist because their constant terms are identities;
the state series and the joint resolvent also converge for \(|z|<1\).
The initial source and the effective operator transform together. A readout
that also acts on \(m_n\) retains its contribution through
\eqref{lib04:memoria-explicita}, in accordance with \eqref{mem:schur-lector}.

Re-entry here updates an incoming record. By contrast, the archiving policy
below retains each complement in a distinct coordinate. Each operation has
its own conservation law. The kernel \(\mathscr K_j\), Green propagators,
and a dimensional force have different types: the first eliminates a state
component, Green propagators solve source equations with support conditions, and
the force adds a dynamical realization and its units.
